% ============================================================
% CHAPTER 5: PROJECT OVERVIEW
% MCDA 5585 / 5586 Major Project Report
% Bhavik Kantilal Bhagat | A00494758 | Saint Mary's University
% ============================================================

\chapter{Project Overview}
\label{chap:project}

% ────────────────────────────────────────────────────────────
\section{Project Background and the Motivating Incident}
\label{sec:project_background}
% ────────────────────────────────────────────────────────────

\subsection{Pre-Co-op Placement Monitoring Posture}

Before my co-op placement began in mid-Mar~2026, the IA-SRE team operated
without a centralized observability stack. Monitoring relied entirely on
an \emph{ad~hoc} model: business teams in Dublin, London, and Halifax
would notice anomalies in fund processing outputs and raise tickets. There
were no CloudWatch alarms, no dashboards aggregating health across the six
platforms, and no mechanism for the SRE team to receive proactive
notification of degrading conditions --- the team was blind to
\emph{saturation} and \emph{error rate} signals until they had already
caused downstream business impact.

\subsection{CAIS-Pricing Silent Outage (Jan--Feb 2026)}
\label{sec:cais_silent_outage}

Between \textbf{Jan~20 and Feb~17, 2026}, the CAIS-Pricing Extract GenAI
pipeline under the \texttt{rpacfs1} platform stopped delivering output for
\textbf{29 consecutive days} without a single alert firing. The root cause
was gradual exhaustion of the external CTI AI model API token, which
controls the async SQS callback that triggers the downstream Lambda
stages; without that callback, three downstream functions dropped to zero
invocations while two upstream functions kept running, making the pipeline
appear healthy. The outage was detected only when business users reported
missing pricing data --- 29~days after onset. The post-mortem
identified the absence of throughput alarms, SQS queue depth monitoring,
and end-to-end output verification as the root causes of the failure going
unnoticed. A sub-task was assigned directly to me as the immediate
remediation and became the starting point for the full monitoring programme.

\subsection{The Meridian Incident (Apr~2026)}

The Meridian platform processes daily statements, reconciliations, and
event notifications across fund administration operations. A surge of
automated system events from Citco Recs overwhelmed
Meridian's OpenSearch consumption layer, causing Kafka consumer lag to
reach approximately \textbf{9~million messages} and halting daily
statement processing for 7.5~hours. The SRE team received no alert; the
Dublin business team escalated over three and a half hours after the
first failure signal. The post-mortem recorded: \textit{``No proactive
monitoring/alerting in place on the SRE side for Meridian event volume
anomalies.''} This became the primary driver for the centralized dashboard project.

\subsection{From Incident to Programme}

Both incidents shared the same root cause: no proactive monitoring. The
post-mortems established two specific gaps --- no Kafka consumer lag alarm,
and no centralized visibility across the IA-IT-SRE Infrastructure --- that
directly motivated the six work streams described in the following
section. Kishor elevated the dashboard project to \textbf{Critical} priority on
May~6, 2026, and my co-op placement pivoted from exploratory dashboard research
to a full-scale observability platform build.

% ────────────────────────────────────────────────────────────
\section{Work Streams Overview}
\label{sec:project_overview}
% ────────────────────────────────────────────────────────────

The two motivating incidents established a clear engineering mandate: build
proactive observability across the entire IA-IT-SRE Infrastructure. Six
work streams were pursued, each directly motivated by the findings and
limitations of the previous. Full technical scope and acceptance criteria
for each deliverable are documented in Chapter~\ref{chap:project_goals}.

\subsection{Work Stream 1: CloudWatch Dashboards and IA Infrastructure Monitoring}

Following the CAIS-Pricing Silent Outage, the objective was to build unified
CloudWatch dashboards starting with Lambda --- the AWS service at the center
of that incident. The scope extended incrementally to SQS queues, EC2 instances,
ECS clusters, EKS nodes, and database services across all six platforms,
establishing the tag-based resource discovery patterns that all later work
built upon.

A fundamental limitation became apparent: CloudWatch dashboards are siloed
per AWS service and do not support cross-service filtering across platforms.
This directly motivated Work Stream~2.
\textit{Key deliverables: IISS-435, IISS-455, IISS-461, IISS-469, IISS-482--486.}

\subsection{Work Stream 2: Grafana AMG Centralized Dashboard}

Following the Meridian Incident, the project was elevated to
\textbf{Critical} priority. The objective was to build a centralized
observability platform using Amazon Managed Grafana (AMG) --- addressing
what CloudWatch dashboards could not: a unified, cross-platform view of the
entire IA-IT-SRE Infrastructure in a single workspace. The broader research
objective was to evaluate whether this AWS-native stack could achieve
Dynatrace-equivalent observability.
\textit{Key deliverable: IISS-487.}

\subsection{Work Stream 3: IA Infrastructure Inventory API}

A critical limitation surfaced during the Grafana build: CloudWatch Metric
Insights SQL does not support pattern-based filtering on dimension values.
To solve this, the Lambda discovery scripts were extended into a full Flask
REST API to discover AWS services dynamically across all platforms, enabling
the Grafana dashboard variables to be populated from live infrastructure
rather than static configuration.
\textit{Key deliverable: IISS-507.}

\subsection{Work Stream 4: RPA Production Support}

In parallel with the engineering work streams, dedicated RPA production
support rotations were carried out as part of the team's Follow the Sun
model. The North America slot (11~AM--7~PM~EST) was covered once every four
weeks, serving as the primary contact for RPA incidents escalated through
the SD Portal, Innovation IT email inbox, IA\_INCIDENTS distribution list,
and Teams channels.

Each shift began with a review of the Confluence handover log from the Manila
shift and a queue scan to identify open items or approaching SLA breaches.
Incidents were triaged, investigated, and either resolved or escalated to
the development team with documented findings.
\textit{Key incident tickets: IISS-706, IISS-538, IISS-741, IISS-876.}

\subsection{Work Stream 5: RPA LLM Agent --- Ask Citco RPA}

This work stream involved a production quality and reliability improvement
programme for the \textbf{Ask Citco RPA} chatbot --- a RAG pipeline using
Amazon Bedrock, OpenSearch Serverless, and 11 Lambda functions. The scope
covered production bug resolution, code quality improvements, structured
logging migration, and X-Ray distributed tracing instrumentation.
\textit{Key deliverable: IISS-824.}

\subsection{Work Stream 6: Citco Enterprise Custom UI Dashboards}

While the Grafana AMG dashboard addressed the primary observability gap,
it remained constrained by fixed panel types and no drag-and-drop layout.
This work stream focused on designing and building a Citco-owned enterprise
monitoring UI by extending the existing \texttt{ia-config-portal} application
(built by Soundarya and Sridhar) with two new dashboard pages.
\textit{Key deliverables: IISS-509, IISS-919. Status: In progress;
implementation continuing in Sep--Dec~2026 extension.}
